\section{The relative Iitaka base and its section lattice}
\label{sec:relative-iitaka}

Our goal is to construct an intermediate base $W$ over the marked-period
quotient $S$ and a divisor on $W$ whose complete divisible section spaces
are those of the original logarithmic base. The rank-one generator on
the generic Iitaka fiber will give this exact comparison.

The coefficient estimates are now available, as are the marked quotient
and both of its positivity comparisons. We will first obtain a nonzero
logarithmic form on the generic period fiber. Only after this
nonnegativity has been established will we form its relative Iitaka
fibration.

We use the marked quotient of \Cref{prop:marked-period}.  Resolve its
compactification and the rational quotient map, obtaining a morphism
\(q:Y_0\to S\) of smooth projective varieties.  Its geometric generic
fiber is connected.  Every source model used below is resolved together
with its reduced boundary, so that the strict transform of the preceding
boundary plus all reduced exceptional divisors is SNC.  A higher source
model always means such a log resolution.  We keep the subscript \(0\) for a fixed such source
model and denote its reduced log boundary by \(D_{Y_0}\); subsequent
models are denoted by \(Y\).  In this section \(E_0\)
denotes the highest Hodge bundle of the descended tensor system on a dense
open subset of \(S\).  The homogeneous vector of degree \(a>0\) is viewed
as a rational tensor
\[
 u_*\in aK_{Y_0}\otimes q^*E_0,
\]
with the logarithmic pole and higher-model properties of
\Cref{prop:logarithmic-vector}.  This notation for a rational tensor
allows the poles prescribed by \(aD_{Y_0}\).

We first record why these source modifications preserve the spaces of
logarithmic pluriforms.  If \(p:Y'\to Y_0\) is a log resolution and
\(D'\) is strict transform plus reduced exceptional divisor, then
\begin{equation}\label{eq:iitaka-log-discrepancy}
 K_{Y'}+D'=p^*(K_{Y_0}+D_{Y_0})+R_p,
 \qquad R_p\ge0\quad\text{and}\quad p_*R_p=0.
\end{equation}
For a blowup of a smooth center of codimension \(c\) contained in
\(e\) boundary branches, the new coefficient in \(R_p\) is \(c-e\ge0\).
For a general log resolution the formula follows from the log canonical
discrepancy inequality for the log smooth pair: each exceptional
coefficient is one plus its discrepancy and is therefore nonnegative.
A rational section of a line bundle on the smooth variety \(Y_0\) is regular if it
has no pole at any prime divisor.  Consequently an effective exceptional
divisor introduces no new rational sections, and, for every \(l\ge0\),
\begin{equation}\label{eq:iitaka-log-sections}
 H^0\bigl(Y',l(K_{Y'}+D')\bigr)
   =H^0\bigl(Y_0,l(K_{Y_0}+D_{Y_0})\bigr).
\end{equation}
These are identifications of rational pluriforms and therefore preserve
their ratios.  The same argument applies over the generic point of a
base.

\subsection{Constructing the relative Iitaka fibration}

Over the generic point of $S$, a nonzero scalar coordinate of $u_*$
will supply a logarithmic pluriform.  The relative Iitaka map then
separates its section ratios, following the classical relative
construction of \cite[\S6, Remark~1]{Iitaka1971}.  On its generic fiber
the nonzero degree-$a$ section space has rank one, so a single generator
factors the entire vector.  The exact comparison with the original
source's complete section spaces will require the valuation argument
that follows this construction.

\begin{lemma}[The relative Iitaka diagram]
\label{lem:relative-iitaka-construction}
For \(q:(Y_0,D_{Y_0})\to S\) and \(u_*\) as above, there are smooth
projective varieties \(Y,W\), a birational morphism \(p:Y\to Y_0\),
and morphisms \(x:Y\to W\) and \(h:W\to S\) with geometrically
connected generic fibers, such that \(h\circ x=q\circ p\).
Put \(D=p_*^{-1}D_{Y_0}+\operatorname{Exc}(p)_{\rm red}\) and
\(k=\kappa(Y_{0,\eta_S},K_{Y_{0,\eta_S}}+D_{Y_0,\eta_S})\).
Then \(k\ge0\), \(\dim W-\dim S=k\), and the geometric generic
fiber of \(x\) has logarithmic Iitaka dimension zero.  Write
\(J=Y_{\eta_W}\) for its generic fiber over \(\C(W)\).

The degree \(a\) used for \(u_*\) itself has a nonzero generator
\(\omega\in H^0(J,a(K_J+D_J))\).  There is a rational tensor
\(u_W\in aK_W\otimes h^*E_0\) for which
\begin{equation}\label{eq:iitaka-vector-factorization}
 u_*=\omega\,x^*u_W.
\end{equation}
\end{lemma}

\begin{figure}[ht]
\centering
\begin{tikzcd}[column sep=large,row sep=large]
Y \arrow[r,"x"] \arrow[d,"p"'] & W \arrow[r,"h"] & S\\
Y_0 \arrow[rru,"q"',bend right=15] &&
\end{tikzcd}
\caption{The relative logarithmic Iitaka diagram over the marked-period
quotient $S$, with $p$ birational.  The generic $x$-fiber $J$ has
$\kappa(J,K_J+D_J)=0$; its degree-$a$ generator $\omega$ gives
$u_*=\omega x^*u_W$.  The generic $h$-fiber has dimension
$k=\kappa(Y_{0,\eta_S},K_{Y_{0,\eta_S}}+D_{Y_0,\eta_S})$.}
\label{fig:intermediate-base}
\end{figure}


\begin{proof}
Over the generic point of \(S\) the coefficient bundle of \(u_*\) is
constant.  At any divisor horizontal over \(S\), its lattice is the
ordinary pullback lattice from the interior of \(S\).  The bounds of
\Cref{prop:logarithmic-vector} consequently give regular logarithmic
coefficients on \(Y_{0,\eta_S}\).  If the bounds were verified after
ramification, regularity descends: the order of a pulled-back ordinary
coefficient is the ramification index times its original order.
Choosing a basis of \((E_0)_{\eta_S}\) and trivializing the base
canonical line therefore gives a nonzero element of
\[
 H^0\bigl(Y_{0,\eta_S},a(K_{Y_{0,\eta_S}}+D_{Y_0,\eta_S})\bigr)
       \otimes_{\C(S)}(E_0)_{\eta_S}.
\]
At least one scalar component is nonzero, proving \(k\ge0\).

Work temporarily over \(K=\C(S)\).  Choose a complete logarithmic
pluricanonical system whose image has dimension \(k\), and resolve its
base ideal.  If its degree is \(b\), its morphism \(\phi:Z\to Z_b\)
satisfies
\begin{equation}\label{eq:iitaka-movable-system}
 b(K_Z+D_Z)\sim\phi^*H+F_b,
 \qquad F_b\ge0,
\end{equation}
where \(H\) is the hyperplane divisor on its projective image.
The extra logarithmic discrepancy in a source resolution is effective,
so the indicated inequality remains valid after further resolution.
Take the Stein factor of \(\phi\), resolve it and the source, and
spread the resulting diagram over a dense open subset of \(S\).
Compactification and resolution give \(Y\xrightarrow{x}W
\xrightarrow{h}S\).  The field \(\C(W)\) is the relative algebraic
closure of the system-image field in \(\C(Y)\); in particular it is
algebraically closed in \(\C(Y)\).  Since \(\C(S)\) is algebraically
closed in \(\C(Y)\), it is also algebraically closed in \(\C(W)\).
In characteristic zero these facts give geometrically connected
generic fibers for both maps.

Here is the section-ratio argument for \(\kappa(J,K_J+D_J)=0\).
The restriction of \(u_*\) still gives a nonzero section in degree
\(a\), exactly as above.  Suppose two independent logarithmic
pluriforms existed on \(J\) in some common degree \(c\).
Their ratio \(r\) would not lie in \(\C(W)\), and would be
transcendental over that field, by its relative algebraic closure.
Spread the two sections to rational sections of
\(c(K_Z+D_Z)\), trivializing the canonical directions from the base
rationally.  Their pole divisors are vertical over \(W_{\eta_S}\).
There are finitely many such prime divisors; the images of all of them
in \(Z_b\) are proper, since \(W_{\eta_S}\to Z_b\) is generically
finite.  Hypersurfaces on \(Z_b\) containing these images, with
sufficient multiplicity, give a section of some \(nH\) whose pullback
cancels all their poles.  Multiply also by the effective section of
\(nF_b\) in \eqref{eq:iitaka-movable-system}.  The two resulting
sections belong to
\(H^0(Z,(c+nb)(K_Z+D_Z))\) and retain ratio \(r\).
Together with the \(k\) independent ratios of the original system,
they give \(k+1\) independent ratios of logarithmic pluriforms.
All finitely many ratios can be placed in one degree by multiplying
numerators and denominators by the other denominators.  This
contradicts the definition of \(k\).

The same conclusion holds for the geometric generic fiber.  In every
fixed degree, proper cohomology in degree zero commutes with field
extension; two sections over an algebraic closure would therefore
already give section-space dimension at least two over \(\C(W)\).
Adjunction and generic smoothness identify the fiber line with
\(K_J+D_J\), and \eqref{eq:iitaka-log-sections} justifies all the log
source replacements.  This proves the claimed geometric Iitaka
dimension without an abundance assumption.

The degree-\(a\) space on \(J\) is now one-dimensional and nonzero.
Choose its generator \(\omega\) and extend it as a rational relative
canonical tensor.  Divide the vector \(u_*\) by it on the generic
fiber.  Its scalar coefficients are in \(\C(W)\), so they define
\(u_W\) and prove \eqref{eq:iitaka-vector-factorization} in the stated
degree \(a\).  The Higgs operator is \(\OO\)-linear; thus its iterates
satisfy the same factorization under differential pullback.  There is
no derivative of \(\omega\) in this assertion.  In particular the
projective direction \([u_W]\) gives precisely the same compact flags
as \([u_*]\) at general corresponding points.
\end{proof}

\subsection{The exact section lattice}

The factorization has so far been obtained over the function field.
To recover the complete logarithmic section spaces on $Y_0$, we must
determine which rational tensors on $W$ multiply by powers of $\omega$
to become regular on the source.  We must test divisors on the fixed
original source even when their images have codimension at least two
on a temporary base. The next lemma extracts precisely those finitely many
base valuations; subsequent source modifications add no new sections.

\begin{lemma}[Extraction of the source valuations]
\label{lem:valuation-preparation}
Let \(x_1:Y_1\to W_1\) be a dominant projective morphism of smooth
projective varieties with geometrically connected generic fiber, and fix
a reduced simple-normal-crossing divisor
\(D_1\) on \(Y_1\).  Write \(J=Y_{1,\eta_{W_1}}\), and suppose that its geometric
logarithmic Iitaka dimension is zero.  Choose, over \(\C(W_1)\),
\[
 0\ne\omega\in H^0\bigl(J,a(K_J+D_J)\bigr)
\]
for some \(a>0\).  Regard \(\omega\) as a rational section of
\(aK_{Y_1/W_1}\).  There is a smooth projective modification of \(W_1\)
with the following property.  On any further smooth base model \(W\),
take a smooth resolution \(p:Y\to Y_1\) of the main component, with
\(x:Y\to W\), and set \(D=p_*^{-1}D_1+\operatorname{Exc}(p)_{\rm red}\).
For a prime \(A\subset W\), put
\begin{equation}\label{eq:iitaka-minimum}
 \begin{aligned}
 t_A&=\min_{I\,\mapsto\,A}
       \frac{a^{-1}\operatorname{ord}_I(\omega)+d_I}{m_I},\\
 m_I&=\operatorname{ord}_I(x^*A),
       \qquad d_I=\operatorname{coeff}_I D,\\
 T_W&=\sum_A t_AA,\qquad L_W=K_W+T_W.
 \end{aligned}
\end{equation}
Here \(I\) runs through prime divisors dominating \(A\), and the order
of \(\omega\) is its order as a section of \(a(K_Y-x^*K_W)\).
The sum defining \(T_W\) is finite.  For every sufficiently divisible
\(l>0\), multiplication gives an isomorphism
\begin{equation}\label{eq:iitaka-exact-sections}
 \begin{split}
 H^0(W,lL_W)&\xrightarrow{\ \simeq\ }
       H^0\bigl(Y_1,l(K_{Y_1}+D_1)\bigr),\\
 \sigma&\longmapsto \omega^{l/a}x^*\sigma.
 \end{split}
\end{equation}
The right side is interpreted on the fixed model \(Y_1\), by rational
identification.  The isomorphisms respect multiplication in a common
Veronese subring and preserve ratios.
\end{lemma}

\begin{proof}
Write \(n=\dim Y_1\), \(w=\dim W_1\), and \(r=n-w\).  We explain the
valuation assertion needed for the modification.  Let \(I_1\) be a prime
of \(Y_1\) whose image has codimension at least two in \(W_1\), and let
\(v=\operatorname{ord}_{I_1}\).  Its restriction to \(\C(W_1)\) is
nontrivial: a local function vanishing on the image has positive order
at \(I_1\).  The nonzero value subgroup is \(m\Z\) for an integer
\(m>0\).  Let \(v_0\) be the normalized restricted valuation and let
\(\kappa(v)\) and \(\kappa(v_0)\) denote their residue fields.

Residues algebraically independent over \(\kappa(v_0)\) lift to elements
algebraically independent over \(\C(W_1)\).  Indeed, in a proposed
polynomial relation among lifts of value zero, divide all coefficients
by a coefficient of smallest valuation.  Reduction then gives a
nonzero polynomial relation among their residues.  It follows that
\[
 \operatorname{trdeg}_{\kappa(v_0)}\kappa(v)\le r,
 \qquad
 \operatorname{trdeg}_{\C}\kappa(v_0)\ge(n-1)-r=w-1.
\]
The reverse inequality is elementary: lifts of independent residues,
together with one element of positive valuation, are algebraically
independent over \(\C\), by examining the term of smallest valuation
in a polynomial relation.  Thus the residue transcendence degree is
exactly \(w-1\).

Choose \(w-1\) such independent residues and rational-function lifts on
\(W_1\).  Resolve the graph of their map to \((\mathbf P^1)^{w-1}\).
The center of \(v_0\) on this proper model has dimension at least
\(w-1\), since its image contains the generic point of the product;
it is proper since the valuation is nontrivial.  It is therefore a
prime divisor.  Resolving further makes the base smooth and does not
contract that center.  This proves that the restriction of \(v\) is a
positive integral multiple of a divisorial valuation of the base.

Only finitely many primes \(I_1\) need this treatment.  Such a prime
lies in the locus where the local fiber dimension of \(x_1\) is at
least \(r+1\).  This is a proper closed subset of \(Y_1\), and a prime
divisor contained in it must be one of its finitely many divisorial
components.  Extract all the restricted valuations simultaneously.
After any further base modification their centers are still divisors.
On a compatible source resolution, the strict transform of each
\(I_1\) survives and dominates its extracted center.  Every prime on
the fixed model \(Y_1\) now either dominates the base or has strict
transform dominating a base prime.

There are finitely many divisors in
\(\operatorname{div}_{aK_{Y/W}}(\omega)+aD\).  Outside the images of
their vertical components, all the numerators in
\eqref{eq:iitaka-minimum} vanish.  This proves finiteness of \(T_W\).
For a nonzero section \(s\) on the right of
\eqref{eq:iitaka-exact-sections}, pullback and
\eqref{eq:iitaka-log-discrepancy} first make \(s\) a logarithmic
pluriform on \((Y,D)\).  If \(a\mid l\), the section space on \(J\)
in degree \(l\) is exactly the line generated by \(\omega^{l/a}\):
it contains that element, and two independent sections would have a
nonconstant ratio, contradicting \(\kappa(J,K_J+D_J)=0\).
Hence, at the function-field level, there is a unique rational
\(l\)-canonical tensor \(\sigma\) on \(W\) such that
\(s=\omega^{l/a}x^*\sigma\).

Fix a base prime \(A\), take a regular local frame of \(lK_W\) at its
generic point, and let \(c\) be the coefficient of \(\sigma\) in that
frame.  For every \(I\) dominating \(A\), regularity of \(s\) says
\begin{equation}\label{eq:iitaka-order-test}
 m_I\operatorname{ord}_A(c)
       +\frac la\operatorname{ord}_I(\omega)+l d_I\ge0.
\end{equation}
These inequalities are equivalent to
\(\operatorname{ord}_A(c)+lt_A\ge0\).  After clearing denominators,
they say precisely that \(\sigma\in H^0(W,lL_W)\).

Conversely, suppose \(\sigma\in H^0(W,lL_W)\), and form the rational
pluriform \(s=\omega^{l/a}x^*\sigma\).  A horizontal prime of the fixed
source \(Y_1\) satisfies the logarithmic regularity test because
\(\omega^{l/a}\) is a logarithmic pluriform on \(J\).  The strict
transform of every other prime of \(Y_1\) dominates a prime \(A\) of
\(W\), by the extraction just performed.  The minimum rule implies
\eqref{eq:iitaka-order-test} there.  The boundary coefficient and the
order of an absolute rational pluriform at a surviving prime agree
with those on \(Y_1\).  Thus \(s\) is regular as a section of
\(l(K_{Y_1}+D_1)\) at every prime of the fixed model.  Smoothness, or
equivalently the normality extension criterion for a line bundle,
extends it across the remaining codimension-two subset.

This argument does not discard the primes contracted by \(x_1\): their
tests are exactly the reason for extracting their restricted valuations.
It also explains why no indefinite extraction of all new source
exceptional divisors is needed.  Once a section is regular on \(Y_1\),
\eqref{eq:iitaka-log-discrepancy} makes its pullback regular on every
higher log source.  Both maps in \eqref{eq:iitaka-exact-sections} are
inverse rational identifications, which proves the last assertions.
\end{proof}

\begin{proposition}[The rank-one base]
\label{prop:rank-one-base}
Take the diagram and degree-\(a\) generator \(\omega\) of
\Cref{lem:relative-iitaka-construction}.  After further smooth
birational preparation, with the reduced source boundary
\(D=p_*^{-1}D_{Y_0}+\operatorname{Exc}(p)_{\rm red}\), the divisor
\(L_W\) defined by \eqref{eq:iitaka-minimum} is big over \(S\).
In every sufficiently divisible degree it satisfies the exact
section comparison
\begin{equation}\label{eq:iitaka-fixed-source}
 H^0(W,lL_W)\simeq
 H^0\bigl(Y_0,l(K_{Y_0}+D_{Y_0})\bigr).
\end{equation}
Moreover, the models can be chosen so that
\begin{equation}\label{eq:iitaka-bundle-formula}
 L_W=K_W+T_W=K_W+B_W+M_W,
 \qquad 0\le B_W\le1,
 \qquad M_W\ \text{nef},
\end{equation}
where \(B_W\) has simple-normal-crossing support and the equality uses
compatible representatives of the rational divisor classes.  If \(I\)
attains the minimum at a prime \(A\), then
\begin{equation}\label{eq:iitaka-discriminant-bound}
 \operatorname{coeff}_A B_W
       \ge 1-\frac{1-d_I}{m_I}\ge0.
\end{equation}

After any further birational preparation of the diagram, the minimum
rule can be recomputed with the higher reduced log boundary.  The
section comparison with the same fixed \((Y_0,D_{Y_0})\) still holds, and
the moduli divisor is the trace of the same b-divisor.  Once that
b-divisor descends to the chosen model, all these higher traces are
pullbacks of \(M_W\).  A further log resolution restores the
simple-normal-crossing condition on \(B_W\).
\end{proposition}

\begin{proof}
\emph{The exact lattice and relative bigness.}
Fix an initial smooth source model \(Y_1\to Y_0\) on which the just
constructed Iitaka map is a morphism.  Apply
\Cref{lem:valuation-preparation} to \(Y_1\), and thereafter use only
models dominating its extracted base.  That lemma, composed with
\eqref{eq:iitaka-log-sections}, proves
\eqref{eq:iitaka-fixed-source} with the original fixed \(Y_0\).
Its order argument also applies over \(\eta_S\).  The logarithmic
systems on \(Y_{\eta_S}\) have image dimension \(k\), and their
ratios descend to systems of \(L_W|_{W_{\eta_S}}\).  After replacing
a degree by a multiple to clear denominators, those descended ratios
still have transcendence degree \(k=\dim W_{\eta_S}\).  Thus \(L_W\)
is big over \(S\).  This remains true when \(k=0\), with the usual
meaning of bigness on a zero-dimensional generic fiber.

\emph{The auxiliary subpair.}
On the current smooth diagram define
\begin{equation}\label{eq:iitaka-auxiliary-pair}
 \Delta=-\frac1a\operatorname{div}_{aK_{Y/W}}(\omega)+x^*T_W.
\end{equation}
Choose compatible rational canonical forms, and write \(\varphi\) for
the rational function expressing \(\omega\) in the corresponding
relative frame.  Then
\begin{equation}\label{eq:iitaka-trivialization}
 K_Y+\Delta+\frac1a\operatorname{div}(\varphi)=x^*L_W.
\end{equation}
At a horizontal prime the inequality \(\Delta\le D\) follows from
the logarithmic regularity of \(\omega|_J\).  At every prime
dominating a fixed base prime \(A\), it is exactly the minimum
inequality \eqref{eq:iitaka-minimum}.  Divisors mapping into
codimension two disappear after shrinking around \(\eta_A\).
It follows that \(\Delta\le D\) near \(x^{-1}(\eta_A)\) for every
\(A\), and also over the generic point of \(W\).  Since the reduced
log smooth pair \((Y,D)\) is log canonical, these inequalities imply
that \((Y,\Delta)\) is sub-log-canonical there.  Negative coefficients
of \(\Delta\) cause no difficulty.

We verify the discrepancy-sheaf rank condition rather than inferring
it from generic log canonicity.  On a resolved geometric generic
fiber \(\rho:J'\to J\), use the higher reduced boundary \(D_{J'}\)
and put
\[
 Z_{J'}=\operatorname{div}_{aK_{J'}}(\omega)+aD_{J'}\ge0.
\]
The crepant transform \(\Delta_{J'}\) of the auxiliary subpair is
\[
 \Delta_{J'}=D_{J'}-\frac1a Z_{J'}.
\]
Moreover \(Z_{J'}\sim a(K_{J'}+D_{J'})\), so
\(\kappa(J',Z_{J'})=0\).  For the discrepancy b-divisor
\(\mathbf A^*\), the coefficient \(-1\) is omitted.  On this
resolution its rounded trace
\(R=\lceil\mathbf A^*_{J'}\rceil\) is effective and is supported
on \(\Supp Z_{J'}\): where the logarithmic zero divisor vanishes,
the auxiliary coefficient is either \(0\) or \(1\), and its
contribution to \(R\) is zero.  Therefore \(R\le N Z_{J'}\) for
some integer \(N\), and
\[
 h^0(J',\OO(R))\le h^0(J',\OO(NZ_{J'}))=1.
\]
If \(Z_{J'}=0\), this simply says \(R=0\).  Sections of the
discrepancy sheaf satisfy in particular the order conditions on this
resolution, so the displayed bound is also an upper bound for its
generic rank.  Conversely, on every higher model sub-log-canonicity
gives discrepancies at least \(-1\); after omitting \(-1\), their
ceilings are nonnegative.  The constant function \(1\) consequently
belongs to the discrepancy sheaf.  We have proved
\begin{equation}\label{eq:iitaka-rank-condition}
 \operatorname{rank}x_*\OO_Y
       \bigl(\lceil\mathbf A^*(Y,\Delta)\rceil\bigr)=1.
\end{equation}

The morphism \(x\) has connected fibers: its Stein factor is finite
birational over the normal variety \(W\), since its geometric generic
fiber is connected.  Thus \(x\) is projective and surjective with connected fibers,
its source and base are normal, \(\Delta\) is rational, the subpair is
sub-log-canonical generically, and
\eqref{eq:iitaka-trivialization} and
\eqref{eq:iitaka-rank-condition} give its remaining lc-trivial
conditions.  Thus it is an lc-trivial fibration in the precise sense of
\cite[Definition~3.2]{FG}.  By \cite[Theorem~3.6]{FG}, applied with
the structural morphism \(W\to\operatorname{Spec}\C\), its moduli
b-divisor is b-nef.  This application requires neither effectivity of
the auxiliary subpair nor abundance or semiampleness of the moduli
part.

\emph{The discriminant and higher models.}
Following \cite[\S3.4]{FG}, define
\[
 c_A=\sup\{c\in\Q:(Y,\Delta+c x^*A)
                  \text{ is sub-log-canonical over }\eta_A\},
 \qquad b_A=1-c_A.
\]
The comparison with \(D\) already proved gives \(c_A\ge0\), hence
\(b_A\le1\).  If \(I\) attains the minimum in
\eqref{eq:iitaka-minimum}, its coefficient in \(\Delta\) is exactly
\(d_I\).  The discrepancy condition at that prime forces
\(d_I+c_A m_I\le1\).  Therefore
\[
 1\ge b_A\ge1-\frac{1-d_I}{m_I}\ge0,
\]
where the last inequality uses \(d_I\in\{0,1\}\) and \(m_I\ge1\).
Put \(B_W=\sum_A b_AA\) and \(M_W=T_W-B_W\), using the
normalization \eqref{eq:iitaka-trivialization}.  This is the moduli
trace for that normalization.

It remains to reconcile b-nefness on a higher model with recomputation
of our minimum rule.  Let \(\beta:W'\to W\) be a smooth birational
base modification, take a compatible source resolution
\(\rho:Y'\to Y\), and use the higher reduced boundary \(D'\).
Write \(L'=K_{W'}+T_{W'}\) for the recomputed minimum divisor and
\(\Delta'\) for \eqref{eq:iitaka-auxiliary-pair} on that diagram.
Let \(\Delta^{\rm cr}\) be the crepant transform, defined by
\(K_{Y'}+\Delta^{\rm cr}=\rho^*(K_Y+\Delta)\).  Comparing
\eqref{eq:iitaka-trivialization} on the two models, with the same
rational trivialization, gives
\begin{equation}\label{eq:iitaka-normalization-change}
 \Delta'=\Delta^{\rm cr}+x'^*Q,
 \qquad Q=L'-\beta^*L_W.
\end{equation}
This identity includes the change of the relative canonical frame;
omitting that change would give an incorrect formula for \(T_{W'}\).

Let \(B^{\rm cr}_{W'}\) be the discriminant of the crepant induced
fibration, as in \cite[\S3.3]{FG}.  At a prime of \(W'\), adding
\(x'^*Q\) decreases its threshold by the local coefficient of
\(Q\).  Thus the recomputed discriminant and moduli divisors satisfy
\begin{equation}\label{eq:iitaka-moduli-compatibility}
 \begin{split}
 B_{W'}&=B^{\rm cr}_{W'}+Q,\\
 M_{W'}&=L'-K_{W'}-B_{W'}
       =\beta^*L_W-K_{W'}-B^{\rm cr}_{W'}.
 \end{split}
\end{equation}
The latter is precisely the old moduli b-divisor's trace.  Hence we
may first dominate a model on which it is nef and descends, and then
make any necessary further modifications; its new trace is its nef
pullback.  Repeating the minimum and threshold argument on each new
diagram still proves \(0\le B_{W'}\le1\).

For completeness, at a base prime not exceptional for \(\beta\) the
minimum does not change.  Indeed, the crepant transform of the old
auxiliary divisor is bounded by \(D'\) over its generic point, by
\eqref{eq:iitaka-log-discrepancy}; this gives the old minimum as a
lower bound for all new candidates.  An old minimizing prime survives
strictly and supplies equality.  The discriminant is likewise
unchanged there.  Thus any new discriminant support is contained in
the exceptional locus together with the transform of the old support.
Resolve that union.  Recomputing the minimum on the resulting model
therefore gives simple-normal-crossing support and keeps the moduli
part nef.  The extraction in \Cref{lem:valuation-preparation} survives
all these operations, so the exact section comparison continues to
refer to the same fixed \(Y_0\).  This proves all the assertions.
\end{proof}

All arguments in this section use the logarithmic vector constructed
above and the published lc-trivial b-nefness theorem.  The logarithmic
Iitaka lower bound of \Cref{thm:log-iitaka} is not used.
